%%
%% Numeric Tower Phase 2: Rationals, MPFR, and Complex Numbers
%%
%% Companion to "Numeric Tower Phase 1: Implementation Proposal"
%%
\documentclass[12pt]{article}
\usepackage[margin=1in]{geometry}
\usepackage{booktabs}
\usepackage{listings}
\usepackage{xcolor}
\usepackage{hyperref}
\usepackage{textcomp}
\usepackage{amssymb}
\usepackage{amsmath}
\usepackage{enumitem}
\usepackage{longtable}

\lstdefinelanguage{Scheme}{
  morekeywords={define,lambda,if,cond,else,let,let*,letrec,begin,
    set!,quote,quasiquote,unquote,unquote-splicing,and,or,not,
    case,do,delay,force,map,for-each,apply,call-with-current-continuation,
    call/cc,require-modules,load,macro,define-global-symbol,
    exact->inexact,inexact->exact,make-rectangular,make-polar,
    real-part,imag-part,magnitude,angle,numerator,denominator,
    rationalize,number->string,string->number,make-mpfr,mpfr-precision},
  sensitive=true,
  morecomment=[l]{;},
  morestring=[b]",
  literate={lambda}{$\lambda$}1,
}

\lstdefinelanguage{Modula3}{
  morekeywords={MODULE,INTERFACE,IMPORT,FROM,PROCEDURE,BEGIN,END,
    VAR,CONST,TYPE,RECORD,OBJECT,METHODS,OVERRIDES,RETURN,
    IF,THEN,ELSE,ELSIF,WHILE,DO,FOR,TO,BY,LOOP,EXIT,REPEAT,UNTIL,
    CASE,OF,WITH,TRY,EXCEPT,FINALLY,RAISE,RAISES,
    NEW,ARRAY,REF,REFANY,SET,BOOLEAN,INTEGER,REAL,LONGREAL,EXTENDED,
    CHAR,TEXT,LONGINT,
    TRUE,FALSE,NIL,BRANDED,REVEAL,EVAL,NOT,AND,OR,IN,DIV,MOD,
    ISTYPE,NARROW,TYPECODE,GENERIC,ADDRESS,BITS,BITSIZE},
  sensitive=true,
  morecomment=[n]{(*}{*)},
  morestring=[b]",
}

\lstset{
  basicstyle=\ttfamily\small,
  breaklines=true,
  frame=single,
  backgroundcolor=\color{gray!5},
  keywordstyle=\bfseries\color{blue!70!black},
  commentstyle=\itshape\color{green!50!black},
  stringstyle=\color{red!60!black},
  showstringspaces=false,
  columns=flexible,
  tabsize=2,
  xleftmargin=2em,
  framexleftmargin=1.5em,
}

\title{\bfseries\sffamily\fontsize{24.88}{30}\selectfont
  Numeric Tower Phase~2:\\Rationals, MPFR, and Complex Numbers}
\author{Mika Nystr\"om \\ {\tt mika@alum.mit.edu}}
\date{March 2026}

\begin{document}
\maketitle
\thispagestyle{empty}

\newpage
\tableofcontents
\newpage

%======================================================================
\section{Overview}
\label{sec:overview}
%======================================================================

Phase~1 of the MScheme numeric tower introduced exact integers
(fixnum and bignum) with overflow-safe arithmetic, exactness
contagion, and full R5RS predicate semantics.  It is complete and
tested: 450/450 numeric tower tests, 184/184 examples, 97/97 R4RS
conformance.

Phase~2 extends the tower with three new numeric types:
\begin{enumerate}[nosep]
\item \textbf{Exact rationals} ({\tt SchemeRational.T}) ---
  pairs of exact integers in lowest terms, so that {\tt (/ 1 3)}
  returns an exact {\tt 1/3} rather than an inexact {\tt 0.333\ldots}.
\item \textbf{Arbitrary-precision inexact} ({\tt SchemeMpfr.T}) ---
  MPFR-backed floating-point values with configurable mantissa
  precision, opt-in via {\tt (make-mpfr \emph{value} \emph{bits})}.
\item \textbf{Complex numbers} ({\tt SchemeComplex.T}) ---
  pairs of arbitrary numeric types, completing the R5RS tower.
\end{enumerate}

\noindent The resulting type hierarchy is:
\begin{lstlisting}[language={},frame=none,backgroundcolor=\color{white},xleftmargin=4em]
SchemeObject.T
+-- Exact
|   +-- SchemeInt.T        (fixnum: cached [-256..256])
|   +-- Mpz.T              (bignum: arbitrary-precision integer)
|   +-- SchemeRational.T   (num/den pair of exact integers)    NEW
+-- Inexact
|   +-- SchemeLongReal.T   (IEEE 754 LONGREAL, 53-bit mantissa)
|   +-- SchemeMpfr.T       (MPFR arbitrary-precision float)    NEW
+-- SchemeComplex.T        (pair of any numeric types)          NEW
\end{lstlisting}

%======================================================================
\section{Build Instructions}
\label{sec:build}
%======================================================================

Build instructions are unchanged from Phase~1 (see the Phase~1
document, Section~2).  Phase~2 adds new modules but does not change
the build procedure.  The new modules require {\tt mpfr} in
addition to {\tt mpz}:

\begin{lstlisting}[language={}]
# In mscheme/src/m3makefile, add:
import("mpfr")
Module("SchemeRational")
Module("SchemeMpfr")
Module("SchemeComplex")
\end{lstlisting}

\noindent The {\tt mpfr} package is already available in the
intel-async tree ({\tt m3utils/mpfr}).  It depends on {\tt mpz}
(and hence on {\tt libgmp} and {\tt libmpfr}).

%======================================================================
\section{Exact Rationals ({\tt SchemeRational})}
\label{sec:rational}
%======================================================================

This is the largest single change and has the widest impact: it
alters the semantics of exact division from returning inexact to
returning exact.

%----------------------------------------------------------------------
\subsection{Type Representation}
%----------------------------------------------------------------------

\begin{lstlisting}[language=Modula3]
TYPE T = BRANDED "SchemeRational" REF
           RECORD num, den: Mpz.T END;
\end{lstlisting}

\noindent Invariants maintained by the constructor:
\begin{enumerate}[nosep]
\item Numerator and denominator are always in lowest terms
  (GCD-normalized).
\item Denominator is always positive.
\item If denominator equals~1, the value is demoted to an exact
  integer ({\tt SchemeInt.T} or {\tt Mpz.T}), analogous to the
  existing bignum-to-fixnum demotion.
\end{enumerate}

The numerator and denominator are stored as {\tt Mpz.T} rather
than {\tt SchemeObject.T} because the normalization algorithm
requires GMP operations ({\tt Mpz.gcd}, {\tt Mpz.tdiv\_q}) and
there is no benefit to distinguishing fixnum-sized components
at the rational level---the {\tt Mpz.T} representation is already
compact for small values.

%----------------------------------------------------------------------
\subsection{Interface}
%----------------------------------------------------------------------

\paragraph{Files.}
\begin{itemize}[nosep]
\item {\tt mscheme/src/SchemeRational.i3} (new)
\item {\tt mscheme/src/SchemeRational.m3} (new)
\end{itemize}

\begin{lstlisting}[language=Modula3]
INTERFACE SchemeRational;
IMPORT SchemeObject, Mpz;

TYPE T = BRANDED "SchemeRational" REF
           RECORD num, den: Mpz.T END;

PROCEDURE New(num, den: Mpz.T): SchemeObject.T;
  (* Normalize and demote to integer if den=1.
     Signals error if den=0. *)

PROCEDURE Add(a, b: T): SchemeObject.T;
PROCEDURE Sub(a, b: T): SchemeObject.T;
PROCEDURE Mul(a, b: T): SchemeObject.T;
PROCEDURE Div(a, b: T): SchemeObject.T;
PROCEDURE Neg(a: T): SchemeObject.T;
PROCEDURE Abs(a: T): SchemeObject.T;

PROCEDURE Compare(a, b: T): INTEGER;
  (* Returns -1, 0, or 1 *)

PROCEDURE IsZero(a: T): BOOLEAN;
PROCEDURE IsPositive(a: T): BOOLEAN;
PROCEDURE IsNegative(a: T): BOOLEAN;

PROCEDURE Numerator(x: T): SchemeObject.T;
  (* Returns exact integer numerator *)
PROCEDURE Denominator(x: T): SchemeObject.T;
  (* Returns exact integer denominator *)

PROCEDURE ToLongReal(x: T): LONGREAL;
  (* num/den as LONGREAL; may lose precision *)

PROCEDURE Format(x: T): TEXT;
  (* Returns "num/den" *)

END SchemeRational.
\end{lstlisting}

%----------------------------------------------------------------------
\subsection{Normalization}
%----------------------------------------------------------------------

The {\tt New} constructor performs GCD normalization:

\begin{lstlisting}[language=Modula3]
PROCEDURE New(num, den: Mpz.T): SchemeObject.T =
  VAR g := Mpz.New();
  BEGIN
    IF Mpz.IsZero(den) THEN
      (* error: division by zero *)
    END;

    (* Ensure denominator is positive *)
    IF Mpz.IsNegative(den) THEN
      Mpz.neg(num, num);
      Mpz.neg(den, den);
    END;

    (* Reduce to lowest terms *)
    Mpz.gcd(g, num, den);
    Mpz.tdiv_q(num, num, g);
    Mpz.tdiv_q(den, den, g);

    (* Demote to integer if den=1 *)
    IF Mpz.IsOne(den) THEN
      RETURN SchemeInt.MpzToScheme(num)
    END;

    RETURN NEW(T, num := num, den := den)
  END New;
\end{lstlisting}

\noindent The demotion rule (denominator~$=$~1~$\Rightarrow$~integer)
is critical: it ensures that {\tt (/ 6 3)} returns the exact integer
{\tt 2}, not the rational {\tt 2/1}.  This parallels the existing
bignum-to-fixnum demotion in {\tt SchemeInt.MpzToScheme}.

%----------------------------------------------------------------------
\subsection{Rational Arithmetic}
%----------------------------------------------------------------------

Standard rules for rational arithmetic, with GCD normalization of
every result:

\begin{align*}
\frac{a}{b} + \frac{c}{d} &= \frac{ad + bc}{bd} \\[4pt]
\frac{a}{b} - \frac{c}{d} &= \frac{ad - bc}{bd} \\[4pt]
\frac{a}{b} \times \frac{c}{d} &= \frac{ac}{bd} \\[4pt]
\frac{a}{b} \div \frac{c}{d} &= \frac{ad}{bc}
\end{align*}

\noindent Each result is passed through {\tt New} which normalizes
and demotes if the denominator is~1.  For example:

\begin{lstlisting}[language=Modula3]
PROCEDURE Add(a, b: T): SchemeObject.T =
  VAR num := Mpz.New();
      den := Mpz.New();
      t   := Mpz.New();
  BEGIN
    Mpz.mul(num, a.num, b.den);   (* a.num * b.den *)
    Mpz.mul(t,   b.num, a.den);   (* b.num * a.den *)
    Mpz.add(num, num, t);         (* numerator = ad + bc *)
    Mpz.mul(den, a.den, b.den);   (* denominator = bd *)
    RETURN New(num, den)
  END Add;
\end{lstlisting}

\paragraph{Optimization note.}
A production implementation should use the standard
``cross-multiply'' optimization that computes
$g = \gcd(b, d)$ first and divides through before
multiplying, to keep intermediate values smaller.  This
is a straightforward enhancement but not essential for
correctness.

%----------------------------------------------------------------------
\subsection{Division Semantics Change}
\label{sec:div-change}
%----------------------------------------------------------------------

This is the single most visible behavioral change in Phase~2.

\paragraph{Phase~1 behavior.}
\begin{lstlisting}[language=Scheme]
(/ 1 3)       ;=> 0.333333333333333  (inexact LONGREAL)
(exact? (/ 1 3))  ;=> #f
\end{lstlisting}

\paragraph{Phase~2 behavior.}
\begin{lstlisting}[language=Scheme]
(/ 1 3)       ;=> 1/3  (exact rational)
(exact? (/ 1 3))  ;=> #t
(/ 1.0 3.0)   ;=> 0.333333333333333  (inexact, unchanged)
(exact->inexact (/ 1 3))  ;=> 0.333333333333333
\end{lstlisting}

\paragraph{Implementation.}
The change is in {\tt SchemeExact.Div} (which currently performs
truncated integer division) and in the {\tt /} dispatch in
{\tt SchemePrimitive.m3}.  When both operands are exact integers
and the division is not exact:

\begin{lstlisting}[language=Modula3]
(* In SchemeExact.Div, replace truncated division: *)
PROCEDURE Div(a, b: SchemeObject.T): SchemeObject.T =
  VAR ma := SchemeInt.ToMpz(a);
      mb := SchemeInt.ToMpz(b);
  BEGIN
    (* Construct rational; New() will normalize and
       demote to integer if remainder is zero *)
    RETURN SchemeRational.New(ma, mb)
  END Div;
\end{lstlisting}

\noindent The existing {\tt quotient}, {\tt remainder}, and
{\tt modulo} primitives are unaffected---they continue to perform
truncated/floored integer division and require integer operands.

\paragraph{Test impact.}
Existing tests that assert {\tt (inexact?\ (/ 1 3))} must be
updated to {\tt (exact?\ (/ 1 3))}.  The Phase~1 test suite
already separates exact-division tests (where the result is an
integer) from inexact-division tests, so the impact is contained.

%----------------------------------------------------------------------
\subsection{Changes to {\tt SchemeExact}}
%----------------------------------------------------------------------

\paragraph{Files.}
\begin{itemize}[nosep]
\item {\tt mscheme/src/SchemeExact.i3} (modify)
\item {\tt mscheme/src/SchemeExact.m3} (modify)
\end{itemize}

\paragraph{{\tt Is()} predicate.}
Currently returns {\tt TRUE} for {\tt REF INTEGER} and {\tt Mpz.T}.
Phase~2 extends this to also accept {\tt SchemeRational.T}:

\begin{lstlisting}[language=Modula3]
PROCEDURE Is(x: SchemeObject.T): BOOLEAN =
  BEGIN
    RETURN x # NIL AND
           (ISTYPE(x, REF INTEGER) OR
            ISTYPE(x, Mpz.T) OR
            ISTYPE(x, SchemeRational.T))
  END Is;
\end{lstlisting}

\paragraph{{\tt IsInteger()} predicate.}
This is where {\tt Is} and {\tt IsInteger} diverge.
{\tt IsInteger} returns {\tt FALSE} for rationals (rationals with
denominator~1 have already been demoted to integers):

\begin{lstlisting}[language=Modula3]
PROCEDURE IsInteger(x: SchemeObject.T): BOOLEAN =
  BEGIN
    RETURN x # NIL AND
           (ISTYPE(x, REF INTEGER) OR
            ISTYPE(x, Mpz.T))
    (* SchemeRational.T is exact but NOT integer *)
  END IsInteger;
\end{lstlisting}

\paragraph{Arithmetic dispatch.}
Each arithmetic procedure in {\tt SchemeExact} must handle four
cross-product cases for the two exact subtypes:

\begin{center}
\begin{tabular}{lcc}
\toprule
& \textbf{Integer} $b$ & \textbf{Rational} $b$ \\
\midrule
\textbf{Integer} $a$ & existing code & promote $a$ to $a/1$ \\
\textbf{Rational} $a$ & promote $b$ to $b/1$ & rational $\times$ rational \\
\bottomrule
\end{tabular}
\end{center}

\noindent The promotion ``integer~$\to$~rational with denominator~1''
is trivial: wrap in {\tt SchemeRational.T} with {\tt den = Mpz.NewInt(1)}.
The {\tt TYPECASE} dispatch is:

\begin{lstlisting}[language=Modula3]
PROCEDURE Add(a, b: SchemeObject.T): SchemeObject.T =
  BEGIN
    TYPECASE a OF
    | SchemeRational.T (ra) =>
        TYPECASE b OF
        | SchemeRational.T (rb) =>
            RETURN SchemeRational.Add(ra, rb)
        ELSE
            RETURN SchemeRational.Add(ra, IntToRat(b))
        END
    ELSE
      TYPECASE b OF
      | SchemeRational.T (rb) =>
          RETURN SchemeRational.Add(IntToRat(a), rb)
      ELSE
        (* existing integer x integer code *)
        ...
      END
    END
  END Add;
\end{lstlisting}

\noindent Where {\tt IntToRat} converts an exact integer to a
rational with denominator~1:
\begin{lstlisting}[language=Modula3]
PROCEDURE IntToRat(x: SchemeObject.T): SchemeRational.T =
  BEGIN
    RETURN NEW(SchemeRational.T,
               num := SchemeInt.ToMpz(x),
               den := Mpz.NewInt(1))
  END IntToRat;
\end{lstlisting}

%----------------------------------------------------------------------
\subsection{Reader: Rational Syntax}
\label{sec:reader-rational}
%----------------------------------------------------------------------

The reader must parse {\tt num/den} syntax.  After detecting a
pure integer token, check whether the next character is~{\tt /}
followed by more digits:

\begin{lstlisting}[language=Modula3]
(* In SchemeInputPort.m3, after integer parse: *)
IF Rd.GetChar(rd) = '/' THEN
  (* Parse denominator *)
  den := ScanInteger(rd);
  IF den = 0 THEN error("zero denominator") END;
  RETURN SchemeRational.New(
           Mpz.NewInt(num), Mpz.NewInt(den))
ELSE
  Rd.UnGetChar(rd);
  RETURN SchemeInt.FromI(num)
END
\end{lstlisting}

\noindent Examples:
\begin{lstlisting}[language=Scheme]
1/3       ;=> SchemeRational: 1/3
-7/4      ;=> SchemeRational: -7/4
6/4       ;=> SchemeRational: 3/2 (normalized)
4/2       ;=> SchemeInt: 2 (demoted)
\end{lstlisting}

%----------------------------------------------------------------------
\subsection{Printer}
%----------------------------------------------------------------------

{\tt SchemeExact.Format} gains a new branch:

\begin{lstlisting}[language=Modula3]
| SchemeRational.T (r) =>
    RETURN Mpz.Format(r.num) & "/" & Mpz.Format(r.den)
\end{lstlisting}

\noindent So {\tt (display 1/3)} prints {\tt 1/3}, and
{\tt (display -7/4)} prints {\tt -7/4}.

%----------------------------------------------------------------------
\subsection{New Primitives}
%----------------------------------------------------------------------

\paragraph{{\tt numerator} and {\tt denominator}.}
R5RS~\S6.2.5 specifies these for all real numbers:

\begin{lstlisting}[language=Scheme]
(numerator 1/3)    ;=> 1
(denominator 1/3)  ;=> 3
(numerator 3)      ;=> 3
(denominator 3)    ;=> 1
(numerator 0.5)    ;=> 1.0
(denominator 0.5)  ;=> 2.0
\end{lstlisting}

\noindent For exact rationals, these extract the components directly.
For exact integers, {\tt numerator} returns the integer and
{\tt denominator} returns~1.  For inexact values, the result is
the inexact equivalent of the exact rational representation of
the value.

\paragraph{{\tt rationalize}.}
R5RS defines {\tt (rationalize $x$ $y$)} to return the simplest
rational within $y$ of~$x$.  The Stern--Brocot tree algorithm
gives the simplest rational in an interval $[x-y, x+y]$:

\begin{lstlisting}[language=Scheme]
(rationalize 0.3 1/10)   ;=> 1/3
(rationalize 3.14 0.01)  ;=> 22/7
\end{lstlisting}

%----------------------------------------------------------------------
\subsection{Predicate Table Update}
%----------------------------------------------------------------------

Table~\ref{tab:predicates} shows the predicate dispatch with
rationals added.

\begin{table}[htbp]
\centering
\begin{tabular}{lccccc}
\toprule
\textbf{Predicate} & {\tt REF INT} & {\tt Mpz.T} &
  {\tt Rational} & {\tt LongReal} & {\tt Complex} \\
\midrule
{\tt number?}   & \checkmark & \checkmark & \checkmark & \checkmark & \checkmark \\
{\tt complex?}  & \checkmark & \checkmark & \checkmark & \checkmark & \checkmark \\
{\tt real?}     & \checkmark & \checkmark & \checkmark & \checkmark & $^*$ \\
{\tt rational?} & \checkmark & \checkmark & \checkmark & if int-val & \\
{\tt integer?}  & \checkmark & \checkmark &            & if int-val & \\
{\tt exact?}    & \checkmark & \checkmark & \checkmark &            & if both \\
{\tt inexact?}  &            &            &            & \checkmark & if either \\
\bottomrule
\end{tabular}
\caption{Predicate dispatch for the full numeric tower.
  ``If int-val'' tests $x = \lfloor x \rfloor$ and finiteness.
  $^*$Complex is real iff imaginary part is zero (but such values
  are demoted, so this case does not arise in practice).
  ``If both/either'' refers to the exactness of the real and
  imaginary parts.}
\label{tab:predicates}
\end{table}


%======================================================================
\section{Arbitrary-Precision Inexact ({\tt SchemeMpfr})}
\label{sec:mpfr}
%======================================================================

The existing inexact representation {\tt SchemeLongReal.T} provides
IEEE~754 64-bit precision (53-bit mantissa, $\sim$16 decimal digits).
{\tt SchemeMpfr.T} wraps MPFR values with configurable precision.

%----------------------------------------------------------------------
\subsection{Design Principles}
%----------------------------------------------------------------------

\begin{enumerate}[nosep]
\item \textbf{Opt-in.}  MPFR values are created explicitly via
  {\tt (make-mpfr \emph{value} \emph{bits})}.  Normal inexact
  arithmetic continues to use {\tt LONGREAL} at hardware speed.
\item \textbf{Contagious.}  Once an operand is {\tt SchemeMpfr.T},
  the result is {\tt SchemeMpfr.T}.  This parallels the existing
  exactness contagion rule.
\item \textbf{Max precision.}  When two {\tt SchemeMpfr.T} values
  interact, the result precision is {\tt MAX(p1, p2)}.
\item \textbf{{\tt LONGREAL} promotion.}  Mixing {\tt LONGREAL}
  with {\tt SchemeMpfr.T} promotes the {\tt LONGREAL} to MPFR
  at the MPFR operand's precision.
\end{enumerate}

%----------------------------------------------------------------------
\subsection{Type Representation}
%----------------------------------------------------------------------

\begin{lstlisting}[language=Modula3]
TYPE T = BRANDED "SchemeMpfr" REF
           RECORD val: Mpfr.T; prec: CARDINAL END;
\end{lstlisting}

\noindent The {\tt prec} field caches the mantissa precision (also
available via the MPFR library, but cached for fast dispatch
decisions).

%----------------------------------------------------------------------
\subsection{Interface}
%----------------------------------------------------------------------

\paragraph{Files.}
\begin{itemize}[nosep]
\item {\tt mscheme/src/SchemeMpfr.i3} (new)
\item {\tt mscheme/src/SchemeMpfr.m3} (new)
\end{itemize}

\begin{lstlisting}[language=Modula3]
INTERFACE SchemeMpfr;
IMPORT SchemeObject, Mpfr;

TYPE T = BRANDED "SchemeMpfr" REF
           RECORD val: Mpfr.T; prec: CARDINAL END;

(* Construction *)
PROCEDURE New(prec: CARDINAL): T;
PROCEDURE FromLR(v: LONGREAL; prec: CARDINAL): T;
PROCEDURE FromExact(v: SchemeObject.T; prec: CARDINAL): T;
  (* Convert any exact number to Mpfr at given precision *)

(* Extraction *)
PROCEDURE ToLR(v: T): LONGREAL;
  (* Truncate to LONGREAL; overflow -> +/-Inf *)
PROCEDURE GetPrec(v: T): CARDINAL;

(* Arithmetic *)
PROCEDURE Add(a, b: T): T;
PROCEDURE Sub(a, b: T): T;
PROCEDURE Mul(a, b: T): T;
PROCEDURE Div(a, b: T): T;
PROCEDURE Neg(a: T): T;
PROCEDURE Abs(a: T): T;
PROCEDURE Sqrt(a: T): T;

(* Transcendentals *)
PROCEDURE Sin(a: T): T;
PROCEDURE Cos(a: T): T;
PROCEDURE Tan(a: T): T;
PROCEDURE Exp(a: T): T;
PROCEDURE Log(a: T): T;

(* Comparison *)
PROCEDURE Compare(a, b: T): INTEGER;

(* Formatting *)
PROCEDURE Format(v: T): TEXT;

END SchemeMpfr.
\end{lstlisting}

\paragraph{Precision propagation.}
All binary operations compute at {\tt MAX(a.prec, b.prec)}:

\begin{lstlisting}[language=Modula3]
PROCEDURE Add(a, b: T): T =
  VAR p := MAX(a.prec, b.prec);
      r := New(p);
  BEGIN
    EVAL Mpfr.Add(r.val, a.val, b.val);
    RETURN r
  END Add;
\end{lstlisting}

%----------------------------------------------------------------------
\subsection{Changes to {\tt SchemeInexact}}
%----------------------------------------------------------------------

\paragraph{Files.}
\begin{itemize}[nosep]
\item {\tt mscheme/src/SchemeInexact.i3} (modify)
\item {\tt mscheme/src/SchemeInexact.m3} (modify)
\end{itemize}

\paragraph{{\tt Is()} predicate.}
Extended to accept both {\tt SchemeLongReal.T} and {\tt SchemeMpfr.T}:

\begin{lstlisting}[language=Modula3]
PROCEDURE Is(x: SchemeObject.T): BOOLEAN =
  BEGIN
    RETURN x # NIL AND
           (ISTYPE(x, SchemeLongReal.T) OR
            ISTYPE(x, SchemeMpfr.T))
  END Is;
\end{lstlisting}

\paragraph{Arithmetic dispatch.}
The pattern parallels the exact-side dispatch.  A helper
{\tt HasMpfr} scans operands for MPFR contagion:

\begin{lstlisting}[language=Modula3]
PROCEDURE Add(a, b: SchemeObject.T): SchemeObject.T =
  BEGIN
    IF ISTYPE(a, SchemeMpfr.T) OR
       ISTYPE(b, SchemeMpfr.T) THEN
      (* Promote both to Mpfr, then operate *)
      RETURN SchemeMpfr.Add(
               ToMpfr(a), ToMpfr(b))
    ELSE
      (* Both are SchemeLongReal.T: existing fast path *)
      WITH va = SchemeLongReal.FromO(a),
           vb = SchemeLongReal.FromO(b) DO
        RETURN SchemeLongReal.FromLR(va + vb)
      END
    END
  END Add;
\end{lstlisting}

\noindent Where {\tt ToMpfr} promotes a {\tt SchemeLongReal.T}
to {\tt SchemeMpfr.T} at the other operand's precision.

%----------------------------------------------------------------------
\subsection{New Primitives}
%----------------------------------------------------------------------

\paragraph{{\tt make-mpfr}.}
Creates an MPFR value from any number:

\begin{lstlisting}[language=Scheme]
(make-mpfr 3.14 128)   ;=> 128-bit Mpfr approximation of 3.14
(make-mpfr 1/3 256)    ;=> 256-bit Mpfr approximation of 1/3
(make-mpfr 42 100)     ;=> 100-bit Mpfr representation of 42
\end{lstlisting}

\noindent Implementation dispatches on the input type:
\begin{itemize}[nosep]
\item {\tt SchemeLongReal.T} $\to$ {\tt SchemeMpfr.FromLR(val, prec)}
\item Exact integer $\to$ {\tt SchemeMpfr.FromExact(val, prec)}
\item {\tt SchemeRational.T} $\to$ compute {\tt num/den} in MPFR
  at the specified precision
\item {\tt SchemeMpfr.T} $\to$ copy with new precision (re-round
  if lowering precision)
\end{itemize}

\paragraph{{\tt mpfr-precision}.}
Returns the mantissa precision in bits:
\begin{lstlisting}[language=Scheme]
(mpfr-precision (make-mpfr 1.0 256))  ;=> 256
\end{lstlisting}

\paragraph{Transcendentals.}
The existing {\tt sin}, {\tt cos}, {\tt exp}, {\tt log}, {\tt sqrt}
primitives gain a dispatch branch: when the argument is
{\tt SchemeMpfr.T}, the corresponding {\tt Mpfr.Sin}, {\tt Mpfr.Cos},
etc.\ is called, giving correctly-rounded results at full precision.
For {\tt SchemeLongReal.T} arguments, the existing {\tt Math}
library functions are used (unchanged).


%======================================================================
\section{Complex Numbers ({\tt SchemeComplex})}
\label{sec:complex}
%======================================================================

%----------------------------------------------------------------------
\subsection{Type Representation}
%----------------------------------------------------------------------

\begin{lstlisting}[language=Modula3]
TYPE T = BRANDED "SchemeComplex" REF
           RECORD re, im: SchemeObject.T END;
\end{lstlisting}

\noindent The {\tt re} and {\tt im} fields are arbitrary Scheme
numbers---integers, rationals, {\tt LONGREAL}, or {\tt SchemeMpfr.T}.
This means that complex rationals (e.g., $\frac{1}{3} + \frac{2}{5}i$)
and complex MPFR values fall out naturally without special-casing.

\paragraph{Demotion invariant.}
If the imaginary part is exactly zero, the constructor demotes
to the real part alone.  This parallels rational demotion
(denominator~1) and bignum demotion (fits in fixnum).

%----------------------------------------------------------------------
\subsection{Interface}
%----------------------------------------------------------------------

\paragraph{Files.}
\begin{itemize}[nosep]
\item {\tt mscheme/src/SchemeComplex.i3} (new)
\item {\tt mscheme/src/SchemeComplex.m3} (new)
\end{itemize}

\begin{lstlisting}[language=Modula3]
INTERFACE SchemeComplex;
IMPORT SchemeObject;
FROM Scheme IMPORT E;

TYPE T = BRANDED "SchemeComplex" REF
           RECORD re, im: SchemeObject.T END;

(* Construction *)
PROCEDURE New(re, im: SchemeObject.T): SchemeObject.T;
  (* Demotes to real if im = 0 *)
PROCEDURE MakeRectangular(re, im: SchemeObject.T):
    SchemeObject.T RAISES {E};
PROCEDURE MakePolar(mag, ang: SchemeObject.T):
    SchemeObject.T RAISES {E};

(* Accessors *)
PROCEDURE RealPart(x: T): SchemeObject.T;
PROCEDURE ImagPart(x: T): SchemeObject.T;
PROCEDURE Magnitude(x: SchemeObject.T):
    SchemeObject.T RAISES {E};
PROCEDURE Angle(x: SchemeObject.T):
    SchemeObject.T RAISES {E};

(* Arithmetic *)
PROCEDURE Add(a, b: T): SchemeObject.T RAISES {E};
PROCEDURE Sub(a, b: T): SchemeObject.T RAISES {E};
PROCEDURE Mul(a, b: T): SchemeObject.T RAISES {E};
PROCEDURE Div(a, b: T): SchemeObject.T RAISES {E};
PROCEDURE Neg(a: T): SchemeObject.T RAISES {E};

(* Formatting *)
PROCEDURE Format(x: T): TEXT;

END SchemeComplex.
\end{lstlisting}

%----------------------------------------------------------------------
\subsection{Complex Arithmetic}
%----------------------------------------------------------------------

Complex arithmetic is defined in terms of the existing
{\tt SchemeNumber} dispatch, so that exactness and precision
contagion compose naturally:

\begin{align*}
(a + bi) + (c + di) &= (a+c) + (b+d)i \\
(a + bi) \times (c + di) &= (ac - bd) + (ad + bc)i \\
\frac{a + bi}{c + di} &= \frac{(ac + bd) + (bc - ad)i}{c^2 + d^2}
\end{align*}

\begin{lstlisting}[language=Modula3]
PROCEDURE Add(a, b: T): SchemeObject.T RAISES {E} =
  BEGIN
    RETURN New(SchemeNumber.Add(a.re, b.re),
               SchemeNumber.Add(a.im, b.im))
  END Add;

PROCEDURE Mul(a, b: T): SchemeObject.T RAISES {E} =
  VAR re := SchemeNumber.Sub(
              SchemeNumber.Mul(a.re, b.re),
              SchemeNumber.Mul(a.im, b.im));
      im := SchemeNumber.Add(
              SchemeNumber.Mul(a.re, b.im),
              SchemeNumber.Mul(a.im, b.re));
  BEGIN
    RETURN New(re, im)
  END Mul;
\end{lstlisting}

\noindent Because {\tt SchemeNumber.Add}, {\tt .Mul}, etc.\ already
handle all combinations of exact/inexact and integer/rational/Mpfr,
the complex module does not need to know about these subtypes.

%----------------------------------------------------------------------
\subsection{Exactness of Complex Numbers}
%----------------------------------------------------------------------

A complex number is exact iff both its real and imaginary parts are
exact:

\begin{lstlisting}[language=Scheme]
(exact? (make-rectangular 1/3 2/5))  ;=> #t
(exact? (make-rectangular 1.0 2.0))  ;=> #f
(exact? (make-rectangular 1 2.0))    ;=> #f
\end{lstlisting}

\noindent The {\tt exact?} and {\tt inexact?} predicates must be
updated to handle {\tt SchemeComplex.T} by checking both components.

%----------------------------------------------------------------------
\subsection{Ordering Restriction}
%----------------------------------------------------------------------

R5RS specifies that {\tt <}, {\tt >}, {\tt <=}, {\tt >=} are
undefined for non-real complex numbers.  The implementation signals
an error:

\begin{lstlisting}[language=Scheme]
(< 3+4i 5+6i)  ;=> ERROR: ordering not defined for complex
(= 3+4i 3+4i)  ;=> #t  (equality IS defined)
\end{lstlisting}

%----------------------------------------------------------------------
\subsection{Changes to {\tt SchemeNumber}}
%----------------------------------------------------------------------

\paragraph{Files.}
\begin{itemize}[nosep]
\item {\tt mscheme/src/SchemeNumber.i3} (modify)
\item {\tt mscheme/src/SchemeNumber.m3} (modify)
\end{itemize}

The {\tt SchemeNumber} dispatch layer gains a third tier for
complex numbers.  The dispatch order becomes:

\begin{enumerate}[nosep]
\item If both operands are exact (integer or rational), use
  {\tt SchemeExact} dispatch.
\item If both operands are real (exact or inexact, but not complex),
  use the existing exact/inexact contagion dispatch.
\item If either operand is complex, promote the non-complex operand
  to a complex with imaginary part zero, then use
  {\tt SchemeComplex} dispatch.
\end{enumerate}

\begin{lstlisting}[language=Modula3]
PROCEDURE Add(a, b: SchemeObject.T):
    SchemeObject.T RAISES {E} =
  BEGIN
    (* Complex contagion *)
    IF ISTYPE(a, SchemeComplex.T) OR
       ISTYPE(b, SchemeComplex.T) THEN
      RETURN SchemeComplex.Add(
               ToComplex(a), ToComplex(b))
    END;

    (* Existing real dispatch *)
    IF SchemeExact.Is(a) THEN
      IF SchemeExact.Is(b) THEN
        RETURN SchemeExact.Add(a, b)
      ELSIF SchemeInexact.Is(b) THEN
        RETURN SchemeInexact.Add(
                 ExactToInexact(a), b)
      END
    ELSIF SchemeInexact.Is(a) THEN
      ...
    END;
    ...
  END Add;
\end{lstlisting}

\noindent Where {\tt ToComplex} wraps a real number as a complex
with imaginary part zero:
\begin{lstlisting}[language=Modula3]
PROCEDURE ToComplex(x: SchemeObject.T):
    SchemeComplex.T =
  BEGIN
    TYPECASE x OF
    | SchemeComplex.T (c) => RETURN c
    ELSE
      RETURN NEW(SchemeComplex.T,
                 re := x, im := SchemeInt.Zero)
    END
  END ToComplex;
\end{lstlisting}

%----------------------------------------------------------------------
\subsection{New Primitives}
%----------------------------------------------------------------------

\begin{lstlisting}[language=Scheme]
(make-rectangular 3 4)   ;=> 3+4i
(make-polar 5 0.9273)    ;=> 3.0+4.0i (approx)
(real-part 3+4i)         ;=> 3
(imag-part 3+4i)         ;=> 4
(magnitude 3+4i)         ;=> 5
(angle 3+4i)             ;=> 0.9272... (atan2)
\end{lstlisting}

\noindent {\tt make-polar} computes {\tt re = mag * cos(ang)} and
{\tt im = mag * sin(ang)}.  When the arguments are exact, the
result is inexact (since trigonometric functions produce inexact
values in general).

%======================================================================
\section{Reader Extensions}
\label{sec:reader}
%======================================================================

All changes are in {\tt SchemeInputPort.m3}.

\begin{table}[htbp]
\centering
\begin{tabular}{lll}
\toprule
\textbf{Syntax} & \textbf{Parsed as} & \textbf{Notes} \\
\midrule
{\tt 1/3} & {\tt SchemeRational} & After integer, check for {\tt /} + digits \\
{\tt -7/4} & {\tt SchemeRational} & Sign applies to numerator \\
{\tt +inf.0} & {\tt SchemeLongReal} (+Inf) & Match as identifier, produce number \\
{\tt -inf.0} & {\tt SchemeLongReal} ($-$Inf) & Same \\
{\tt +nan.0} & {\tt SchemeLongReal} (NaN) & Same \\
{\tt 3+4i} & {\tt SchemeComplex} & After number, check for {\tt +}/{\tt -} number {\tt i} \\
{\tt 3-4i} & {\tt SchemeComplex} & Same \\
{\tt +2i} & {\tt SchemeComplex} (0+2i) & Special case: real part is 0 \\
{\tt -3i} & {\tt SchemeComplex} (0$-$3i) & Same \\
{\tt +i} & {\tt SchemeComplex} (0+1i) & Special case: coefficient is 1 \\
{\tt -i} & {\tt SchemeComplex} (0$-$1i) & Same \\
\bottomrule
\end{tabular}
\caption{Reader extensions for Phase~2.}
\label{tab:reader}
\end{table}

\paragraph{Parsing order.}
The number parser proceeds in stages:
\begin{enumerate}[nosep]
\item Scan optional sign and digits.
\item If {\tt /} follows, parse denominator $\to$ rational.
\item If {\tt .} or exponent marker follows, parse as float.
\item Otherwise, parse as integer.
\item After obtaining the first number, if {\tt +} or {\tt -}
  follows (and is not a delimiter), parse a second number and
  expect~{\tt i} $\to$ complex.
\item If the token is exactly {\tt +i} or {\tt -i}, return
  complex with coefficient~1.
\end{enumerate}

\paragraph{Special float literals.}
The tokens {\tt +inf.0}, {\tt -inf.0}, and {\tt +nan.0} are
recognized in the identifier scanner (since they begin with
{\tt +} or {\tt -}, which are valid identifier-start characters
in R7RS).  When matched, they return the corresponding IEEE~754
{\tt LONGREAL} value instead of a symbol.

%======================================================================
\section{Scheme--Modula-3 Interop}
\label{sec:interop}
%======================================================================

When Scheme code calls Modula-3 procedures via {\tt modula3scheme}
stubs (generated by {\tt sstubgen}), numeric types must convert at
the boundary.  Phase~2 adds new source types (rational, Mpfr,
complex) that must map to Modula-3 formal parameter types.

%----------------------------------------------------------------------
\subsection{Scheme $\to$ Modula-3 ({\tt ToModula} Direction)}
%----------------------------------------------------------------------

Table~\ref{tab:to-modula} shows the conversion matrix.

\begin{table}[htbp]
\centering
\small
\begin{tabular}{lccccccc}
\toprule
\textbf{M3 formal} & {\tt Int} & {\tt Mpz} &
  {\tt Rational} & {\tt LongReal} & {\tt Mpfr} & {\tt Complex} & \textbf{Non-num} \\
\midrule
{\tt INTEGER}  & direct & range & error & round & error & error & error \\
{\tt CARDINAL} & $\geq 0$ & range & error & round & error & error & error \\
{\tt LONGREAL} & float & {\tt get\_d} & $n/d$ & direct & {\tt GetLR} & error & error \\
{\tt REAL}     & float & {\tt get\_d} & $n/d$ & trunc & {\tt GetLR} & error & error \\
{\tt EXTENDED} & float & {\tt get\_d} & $n/d$ & trunc & {\tt GetLR} & error & error \\
{\tt Mpz.T}    & {\tt NewInt} & direct & error & error & error & error & error \\
{\tt Mpfr.T}   & {\tt SetInt} & {\tt set\_d} & $n/d$ & {\tt SetLR} & direct & error & error \\
{\tt TEXT}     & error & error & error & error & error & error & {\tt ToText} \\
{\tt BOOLEAN}  & \multicolumn{7}{c}{{\tt TruthO}: all values yield {\tt \#t} except {\tt \#f}} \\
{\tt CHAR}     & error & error & error & error & error & error & char$^*$ \\
{\tt REFANY}   & pass & pass & pass & pass & pass & pass & pass \\
\bottomrule
\end{tabular}
\caption{Scheme $\to$ Modula-3 type conversions.
  ``Range'' means range-check via {\tt Mpz.get\_si};
  ``$n/d$'' means compute numerator/denominator as the target type;
  ``error'' means the conversion is not supported and raises a
  Scheme exception.
  ``Non-num'' covers strings, booleans, characters, and symbols.
  $^*$Only Scheme character objects ({\tt \#\textbackslash A}) are accepted.}
\label{tab:to-modula}
\end{table}

\paragraph{Key rules.}
\begin{itemize}[nosep]
\item \textbf{Mpfr.T $\to$ LONGREAL}: truncate precision;
  overflow yields $\pm\infty$ (no exception).
\item \textbf{Rational $\to$ LONGREAL}: compute {\tt num/den} as
  double; exact if representable.
\item \textbf{Complex $\to$ any real type}: error.  Cannot silently
  discard the imaginary part.
\item \textbf{Rational $\to$ INTEGER} or {\tt Mpz.T}: error.
  A rational is not an integer (rationals with denominator~1
  have already been demoted).
\item \textbf{EXTENDED}: same conversion path as {\tt REAL}
  (via {\tt SchemeLongReal.FromO}, then {\tt FLOAT} to target).
  \emph{Note:} CM3 master is transitioning {\tt EXTENDED} to
  IEEE~754 binary128 (128-bit quadruple precision).  When that
  change is adopted, {\tt ToScheme\_EXTENDED} will demote to
  {\tt LONGREAL} with precision loss.  A future enhancement could
  use {\tt SchemeMpfr.T} to represent {\tt EXTENDED} losslessly.
\item \textbf{TEXT}: only accepts {\tt SchemeString.T};
  all other types raise an exception.
\item \textbf{BOOLEAN}: {\tt TruthO} returns {\tt TRUE} for every
  value except {\tt \#f} (standard Scheme truthiness).
\item \textbf{CHAR}: only accepts {\tt SchemeChar.T} objects
  (e.g., {\tt \#\textbackslash A}); all other types raise an exception.
\end{itemize}

%----------------------------------------------------------------------
\subsection{Modula-3 $\to$ Scheme ({\tt ToScheme} Direction)}
%----------------------------------------------------------------------

\begin{table}[htbp]
\centering
\begin{tabular}{ll}
\toprule
\textbf{M3 return type} & \textbf{Scheme result} \\
\midrule
{\tt INTEGER}  & {\tt SchemeInt.T} (fixnum, cached) \\
{\tt CARDINAL} & {\tt SchemeInt.T} \\
{\tt LONGREAL} & {\tt SchemeLongReal.T} \\
{\tt REAL}     & {\tt SchemeLongReal.T} (promoted to {\tt LONGREAL}) \\
{\tt EXTENDED} & {\tt SchemeLongReal.T} (promoted to {\tt LONGREAL}) \\
{\tt Mpz.T}    & {\tt SchemeInt.T} or {\tt Mpz.T} \\
{\tt Mpfr.T}   & {\tt SchemeMpfr.T} (preserves full precision) \\
{\tt TEXT}     & {\tt SchemeString.T} \\
{\tt BOOLEAN}  & {\tt SchemeBoolean.T} \\
{\tt CHAR}     & {\tt SchemeChar.T} \\
{\tt REFANY}   & pass-through \\
\bottomrule
\end{tabular}
\caption{Modula-3 $\to$ Scheme type conversions.}
\label{tab:to-scheme}
\end{table}

\noindent The key change is that {\tt Mpfr.T} return values stay as
{\tt SchemeMpfr.T} rather than being converted to {\tt LONGREAL}.
This preserves the full precision of MPFR computations across the
Scheme--M3 boundary.

%----------------------------------------------------------------------
\subsection{{\tt sstubgen} Changes}
%----------------------------------------------------------------------

\paragraph{Files.}
\begin{itemize}[nosep]
\item {\tt modula3scheme/src/SchemeModula3Types.i3} (modify)
\item {\tt modula3scheme/src/SchemeModula3Types.m3} (modify)
\end{itemize}

The {\tt Type} enumeration gains a new member:
\begin{lstlisting}[language=Modula3]
TYPE
  Type = { ..., T_Mpz_T, T_Mpfr_T };  (* add T_Mpfr_T *)
\end{lstlisting}

\noindent New conversion procedures:

\begin{lstlisting}[language=Modula3]
PROCEDURE ToModula_Mpfr_T(x: SchemeObject.T):
    Mpfr.T RAISES {Scheme.E} =
  BEGIN
    TYPECASE x OF
    | SchemeMpfr.T (m) =>
        RETURN m.val
    | SchemeLongReal.T (lr) =>
        VAR r := Mpfr.New(53); BEGIN
          EVAL Mpfr.SetLR(r, lr.val);
          RETURN r
        END
    | REF INTEGER (ri) =>
        VAR r := Mpfr.New(64); BEGIN
          EVAL Mpfr.SetInt(r, ri^);
          RETURN r
        END
    | Mpz.T (m) =>
        VAR r := Mpfr.New(Mpz.SizeInBits(m)); BEGIN
          EVAL Mpfr.SetLR(r, Mpz.get_d(m));
          RETURN r
        END
    | SchemeRational.T (q) =>
        VAR rn := Mpfr.New(128);
            rd := Mpfr.New(128); BEGIN
          EVAL Mpfr.SetLR(rn, Mpz.get_d(q.num));
          EVAL Mpfr.SetLR(rd, Mpz.get_d(q.den));
          EVAL Mpfr.Div(rn, rn, rd);
          RETURN rn
        END
    ELSE
      RAISE Scheme.E("expected a number for Mpfr.T")
    END
  END ToModula_Mpfr_T;

PROCEDURE ToScheme_Mpfr_T(m: Mpfr.T):
    SchemeObject.T =
  BEGIN
    RETURN NEW(SchemeMpfr.T,
               val := m,
               prec := Mpfr.GetPrec(m))
  END ToScheme_Mpfr_T;
\end{lstlisting}

\paragraph{Generated stubs.}
Currently, per-package stubs for {\tt Mpfr.T} parameters use
{\tt NARROW(x, Mpfr.T)}, which only accepts actual {\tt Mpfr.T}
objects.  After this change, {\tt sstubgen} generates
{\tt SchemeModula3Types.ToModula\_Mpfr\_T(x)} instead, accepting
any numeric type.  This requires adding {\tt Mpfr.T} as a
concrete-mode base type in {\tt sstubgen}, parallel to how
{\tt Mpz.T} is already handled.


%======================================================================
\section{Contagion Rules}
\label{sec:contagion}
%======================================================================

Phase~2 introduces two new contagion dimensions in addition to
the existing exactness contagion.  All three compose orthogonally.

%----------------------------------------------------------------------
\subsection{Exactness Contagion (Existing)}
%----------------------------------------------------------------------

\begin{center}
\begin{tabular}{lccc}
\toprule
& \textbf{Exact} & \textbf{Inexact} \\
\midrule
\textbf{Exact}   & exact   & inexact \\
\textbf{Inexact} & inexact & inexact \\
\bottomrule
\end{tabular}
\end{center}

\noindent Unchanged from Phase~1 and R5RS.

%----------------------------------------------------------------------
\subsection{Precision Contagion (New)}
%----------------------------------------------------------------------

\begin{center}
\begin{tabular}{lccc}
\toprule
& \textbf{LONGREAL} & \textbf{Mpfr($p$)} \\
\midrule
\textbf{LONGREAL}   & LONGREAL   & Mpfr($p$) \\
\textbf{Mpfr($p'$)} & Mpfr($p'$) & Mpfr(max($p$,$p'$)) \\
\bottomrule
\end{tabular}
\end{center}

\noindent When an exact value meets an {\tt SchemeMpfr.T}, the
exact value is converted to MPFR at the MPFR operand's precision.

%----------------------------------------------------------------------
\subsection{Complex Contagion (New)}
%----------------------------------------------------------------------

\begin{center}
\begin{tabular}{lccc}
\toprule
& \textbf{Real} & \textbf{Complex} \\
\midrule
\textbf{Real}    & real    & complex \\
\textbf{Complex} & complex & complex \\
\bottomrule
\end{tabular}
\end{center}

\noindent A real operand is promoted to complex with imaginary
part zero.

%----------------------------------------------------------------------
\subsection{Combined Contagion}
%----------------------------------------------------------------------

The three contagion rules compose.  Consider:

\begin{lstlisting}[language=Scheme]
(+ (make-rectangular 1/3 2/5)   ; exact complex rational
   (make-mpfr 1.0 128))         ; inexact Mpfr
\end{lstlisting}

\noindent This involves:
\begin{enumerate}[nosep]
\item Complex contagion: promote the Mpfr scalar to complex
  (imaginary part = exact zero).
\item Exactness contagion: exact rational parts are converted
  to inexact.
\item Precision contagion: the conversion yields Mpfr(128), not
  LONGREAL.
\end{enumerate}

\noindent Result: an inexact complex with both parts as
Mpfr(128) values.


%======================================================================
\section{Implementation Order}
\label{sec:order}
%======================================================================

The three new types are implemented in dependency order, each step
self-contained and independently testable.

%----------------------------------------------------------------------
\subsection{Step~1: Exact Rationals}
%----------------------------------------------------------------------

This is the largest step and has the widest behavioral impact
(division semantics change).

\begin{enumerate}[nosep]
\item Create {\tt SchemeRational.i3} and {\tt .m3}.
\item Modify {\tt SchemeExact} to accept rationals in {\tt Is()},
  extend arithmetic dispatch to handle rational~$\times$~integer.
\item Change {\tt SchemeExact.Div} to return rationals.
\item Update reader for {\tt num/den} syntax.
\item Add {\tt numerator}, {\tt denominator}, {\tt rationalize}
  primitives.
\item Update test suite; fix tests affected by division change.
\item Add {\tt SchemeRational} $\to$ {\tt LONGREAL}/{\tt Mpfr.T}
  conversion in {\tt SchemeModula3Types}.
\end{enumerate}

%----------------------------------------------------------------------
\subsection{Step~2: SchemeMpfr}
%----------------------------------------------------------------------

Opt-in and minimally disruptive.

\begin{enumerate}[nosep]
\item Create {\tt SchemeMpfr.i3} and {\tt .m3}.
\item Modify {\tt SchemeInexact} to accept {\tt SchemeMpfr.T},
  add Mpfr arithmetic dispatch.
\item Add {\tt make-mpfr} and {\tt mpfr-precision} primitives.
\item Add Mpfr dispatch to transcendentals.
\item Add {\tt Mpfr.T} as concrete base type in {\tt sstubgen}.
\item Add precision tests to test suite.
\end{enumerate}

%----------------------------------------------------------------------
\subsection{Step~3: Complex Numbers}
%----------------------------------------------------------------------

Depends on rationals for full generality (complex rationals).

\begin{enumerate}[nosep]
\item Create {\tt SchemeComplex.i3} and {\tt .m3}.
\item Modify {\tt SchemeNumber} for complex dispatch.
\item Add {\tt make-rectangular}, {\tt make-polar}, {\tt real-part},
  {\tt imag-part}, {\tt magnitude}, {\tt angle} primitives.
\item Update {\tt exact?}/{\tt inexact?} for complex values.
\item Add {\tt a+bi} reader syntax.
\item Add complex arithmetic tests.
\end{enumerate}

%----------------------------------------------------------------------
\subsection{Step~4: Reader Cleanup}
%----------------------------------------------------------------------

Independent of the above; can be done at any time.

\begin{enumerate}[nosep]
\item {\tt +inf.0}, {\tt -inf.0}, {\tt +nan.0} reader support.
\item Printer round-trip for special float values.
\end{enumerate}


%======================================================================
\section{Files to Create and Modify}
\label{sec:files}
%======================================================================

Table~\ref{tab:files} lists all files affected by Phase~2.

\begin{table}[htbp]
\centering
\small
\begin{tabular}{llp{5.5cm}}
\toprule
\textbf{File} & \textbf{Action} & \textbf{Change} \\
\midrule
{\tt SchemeRational.i3} & New & Exact rational interface \\
{\tt SchemeRational.m3} & New & GCD normalization, rational arithmetic \\
{\tt SchemeMpfr.i3} & New & MPFR wrapper interface \\
{\tt SchemeMpfr.m3} & New & Precision-tracked MPFR operations \\
{\tt SchemeComplex.i3} & New & Complex number interface \\
{\tt SchemeComplex.m3} & New & Complex arithmetic, demotion \\
\midrule
{\tt SchemeExact.i3} & Modify & {\tt Is()} accepts rationals;
  {\tt IsInteger()} diverges \\
{\tt SchemeExact.m3} & Modify & Rational~$\times$~integer dispatch;
  division returns rationals \\
{\tt SchemeInexact.i3} & Modify & {\tt Is()} accepts Mpfr \\
{\tt SchemeInexact.m3} & Modify & Mpfr arithmetic dispatch;
  {\tt ToMpfr} promotion \\
{\tt SchemeNumber.i3} & Modify & {\tt Is()} accepts complex \\
{\tt SchemeNumber.m3} & Modify & Complex dispatch tier;
  {\tt ToComplex} promotion \\
{\tt SchemePrimitive.m3} & Modify & New primitives; division
  change; transcendental dispatch \\
{\tt SchemeInputPort.m3} & Modify & Rational, complex, inf/nan
  reader syntax \\
{\tt mscheme/src/m3makefile} & Modify & Add new modules,
  {\tt import("mpfr")} \\
\midrule
{\tt SchemeModula3Types.i3} & Modify & {\tt T\_Mpfr\_T} type;
  Mpfr/Rational conversions \\
{\tt SchemeModula3Types.m3} & Modify & {\tt ToModula\_Mpfr\_T},
  {\tt ToScheme\_Mpfr\_T} \\
\midrule
{\tt numeric-tower-phase2.tex} & New & This document \\
{\tt test-numeric-tower.scm} & Modify & Add Phase~2 test sections \\
\bottomrule
\end{tabular}
\caption{Files created or modified in Phase~2.}
\label{tab:files}
\end{table}

%======================================================================
\section{Testing}
\label{sec:testing}
%======================================================================

%----------------------------------------------------------------------
\subsection{Existing Tests}
%----------------------------------------------------------------------

All 450 existing numeric tower tests must continue to pass, with
one systematic change: tests that assert {\tt (inexact? (/ $a$ $b$))}
for exact integer operands with non-zero remainder must be updated
to assert {\tt (exact? (/ $a$ $b$))} and check the rational result.

The 184 example tests and 97 R4RS conformance tests should be
unaffected (they do not depend on the inexactness of exact division).

%----------------------------------------------------------------------
\subsection{New Test Sections}
%----------------------------------------------------------------------

The following sections are added to {\tt test-numeric-tower.scm}:

\paragraph{Exact rationals.}
\begin{lstlisting}[language=Scheme]
;; Construction and normalization
(test (/ 1 3)       '1/3)      ; exact rational
(test (/ 6 4)       '3/2)      ; normalized
(test (/ 4 2)       2)         ; demoted to integer
(test (/ -7 4)      '-7/4)     ; sign on numerator
(test (/ 7 -4)      '-7/4)     ; denominator always positive

;; Arithmetic
(test (+ 1/3 1/6)   '1/2)
(test (- 1/2 1/3)   '1/6)
(test (* 2/3 3/4)   '1/2)
(test (/ 2/3 4/5)   '5/6)

;; Mixed integer/rational
(test (+ 1 1/3)     '4/3)
(test (* 3 1/3)     1)         ; demoted to integer

;; Predicates
(test (exact? 1/3)    #t)
(test (rational? 1/3) #t)
(test (integer? 1/3)  #f)
(test (real? 1/3)     #t)

;; Primitives
(test (numerator 1/3)   1)
(test (denominator 1/3) 3)
(test (numerator 3)     3)
(test (denominator 3)   1)
\end{lstlisting}

\paragraph{Arbitrary-precision inexact.}
\begin{lstlisting}[language=Scheme]
;; Construction
(test (mpfr-precision (make-mpfr 3.14 128)) 128)

;; Precision contagion
(test (mpfr-precision
        (+ (make-mpfr 1.0 128)
           (make-mpfr 2.0 256)))
      256)

;; Exactness: make-mpfr produces inexact
(test (inexact? (make-mpfr 42 100)) #t)

;; Transcendentals at high precision
(let ((pi (make-mpfr 0 256)))
  ;; (* 4 (atan 1)) should give pi to 256-bit precision
  ...)
\end{lstlisting}

\paragraph{Complex numbers.}
\begin{lstlisting}[language=Scheme]
;; Construction
(test (make-rectangular 3 4) '3+4i)
(test (make-rectangular 3 0) 3)    ; demoted

;; Arithmetic
(test (+ 1+2i 3+4i) '4+6i)
(test (* 1+2i 3+4i) '-5+10i)

;; Accessors
(test (real-part 3+4i) 3)
(test (imag-part 3+4i) 4)
(test (magnitude 3+4i) 5)

;; Predicates
(test (complex? 3+4i) #t)
(test (number? 3+4i)  #t)
(test (real? 3+4i)    #f)
(test (exact? (make-rectangular 1/3 2/5)) #t)

;; Ordering error
(test-error (< 3+4i 5+6i))

;; Complex rationals
(test (+ (make-rectangular 1/3 1/4)
         (make-rectangular 1/6 3/4))
      (make-rectangular 1/2 1))
\end{lstlisting}

%----------------------------------------------------------------------
\subsection{Integration Testing}
%----------------------------------------------------------------------

After Phase~2, the downstream packages {\tt numtest}, {\tt cspc},
and {\tt photopic} must be rebuilt and their test suites rerun.  The
primary risk is the division semantics change: any Scheme code that
expects {\tt (/ $a$ $b$)} to return a {\tt LONGREAL} when both
operands are exact integers will now receive a rational.  If such a
value is then passed to a Modula-3 procedure expecting {\tt LONGREAL},
the existing {\tt ToModula\_LONGREAL} will handle the conversion
(via the rational~$\to$~double path in Table~\ref{tab:to-modula}).

%======================================================================
\section{What Phase~2 Does Not Change}
\label{sec:not-changed}
%======================================================================

\begin{itemize}
\item \textbf{No compiler changes.}  The MScheme compiler
  ({\tt compiler.scm}) continues to emit {\tt SchemeInt}-based
  code for integer constants and {\tt SchemeLongReal}-based code
  for float constants.  Rational, Mpfr, and complex literals in
  compiled code are handled at runtime by the library, not by
  special compiler support.  Compiler optimizations for the new
  types are a future phase.

\item \textbf{No GMP-free fallback for MPFR.}  Unlike the GMP-free
  fallback discussed in Phase~1 (Section~10), there is no pure-M3
  fallback for MPFR.  If MPFR is not available, {\tt make-mpfr}
  signals an error and the system operates with {\tt LONGREAL}
  only.

\item \textbf{No interval arithmetic.}  MPFR supports directed
  rounding modes that could be used for interval arithmetic, but
  this is not exposed in Phase~2.

\item \textbf{No {\tt exact->inexact} for rationals via continued
  fractions.}  The {\tt inexact->exact} conversion of a
  {\tt LONGREAL} to a rational uses the standard IEEE~754 bit
  decomposition (the {\tt LONGREAL} value is exactly
  $m \times 2^e$ for integer $m$ and $e$).  Continued-fraction
  approximation (which would give ``nicer'' rationals) is not
  implemented.
\end{itemize}

\end{document}
